% !TEX root = ./Basilisk-ThrusterForces-20160627.tex


\section{Module Functions}
This module has the following functions:
\begin{itemize}
	\item \textbf{Read in the thruster and spacecraft configuration}: When reset the module reads in the thruster locations, force directions and maximum thrust values, as well as the spacecraft center of mass.  The control axes $\hat{\bm c}_{j}$ are normalized, and a warning is logged if any two control axes are not orthogonal.
	\item \textbf{Compute the thruster force solution}: With each update call the module reduces the commanded torque $\bm L_{r}$ onto the control axes and computes the minimum norm thruster force solution $\bm F$ that produces $\bar{\bm L}_{r}$.  If {\tt use2ndLoop} is set or off-pulsing is used ({\tt thrForceSign} $= -1$), the 2-stage solution is used to enforce the desired force sign.
	\item \textbf{Evaluate the torque angle error}: The angle $\theta_{\bm\tau/\bar{\bm L}_{r}}$ between the saturated thruster torque $\bm\tau$ and the reduced torque $\bar{\bm L}_{r}$ is evaluated in body frame components, and is stored in {\tt outTorqAngErr}.
	\item \textbf{Scale the thruster forces if saturated}: If this angle is larger than {\tt angErrThresh}, the thruster forces are scaled such that no thruster exceeds its maximum thrust.
	\item \textbf{Output the thruster force commands}: The resulting thruster forces are written to the thruster force output message.
\end{itemize}

\section{Module Assumptions and Limitations}
The module assumes all thrusters in the configuration message are operating and available.  It also assumes the thruster configuration and the spacecraft center of mass don't change between module resets.  Up to three control axes are supported, and these must be orthogonal for the reduced torque $\bar{\bm L}_{r}$ to be valid.
